

\begin{abstract}
Designing RNA sequences that reliably adopt specified three-dimensional structures while maintaining thermodynamic stability remains a central challenge for synthetic biology and therapeutics. Current inverse-folding approaches optimize sequence recovery or a single structural proxy, ignoring two RNA-specific properties that make multi-objective optimization unstable: \emph{heterogeneous noise} across physical proxies (deterministic 2D folding, stochastic 3D prediction, GC-confounded free energy) and \emph{sequence-structure degeneracy} that enables compositional reward hacking via GC enrichment.
We introduce \textbf{RiboPO}, a preference-optimization framework that formalizes RNA inverse folding as preference learning under multiple noisy physical proxies. RiboPO constructs preference pairs as $\varepsilon$-Pareto-dominance relations on standardized per-metric features, gates winners with a structural quality threshold, and trains a frozen-reference DPO policy under a decreasing-margin curriculum. The resulting algorithm admits a trust-region characterization (Theorem~\ref{thm:trust_region}) and an off-policy bias bound (Theorem~\ref{thm:off_policy}) that together explain why three RNA-specific design choices --- variability-aware margins, structural quality gates, and frozen reference --- are each individually necessary; ablations confirm each is load-bearing. We further give two practical extensions: (i) a \emph{thermodynamic-surplus} preference filter that re-fits MFE on length and GC and drops GC-driven pairs, removing the compositional confound at the source; and (ii) a Stage-2 \emph{weight-conditioned policy} via FiLM that adds 1024 parameters and lets users traverse the Pareto front at inference time without retraining.
On the DAS benchmark (98 structures), RiboPO improves EternaFold scMCC ($0.61\!\to\!0.69$, $+13\%$, paired Wilcoxon $p{=}1.2\!\times\!10^{-3}$), Vienna MFE ($-30.4\!\to\!-34.0$ kcal/mol, $-11.8\%$, $p{=}7.0\!\times\!10^{-6}$), target-structure probability $P(S_0)$ (\textbf{$+687\%$}, $p{=}1.6\!\times\!10^{-5}$), and melting temperature ($+1.2\,^\circ$C); designability rates rise by $13\!-\!16$ percentage points. A GC-controlled regression attributes \textbf{$69\%$ of the MFE improvement to GC-independent thermodynamic gain} ($p{=}1.3\!\times\!10^{-7}$); the thermodynamic-surplus filter, applied prospectively, drops $\approx 24\%$ of preference pairs as GC-driven, providing a constructive (not merely post-hoc) defense. Under a joint pLDDT-and-RMSD success criterion, RiboPO pass@1 ($0.258$) exceeds gRNAde pass@64 ($0.154$), a \textbf{$64\times$ sample-efficiency gain}. The framework transfers to a second base model (RiFold) under conservative hyperparameters. Across 21 SSTT metrics, $20$ improve and $17$ are statistically significant; the lone non-improver is sequence recovery, an explicit Pareto trade-off rather than functional degradation --- INF for non-canonical contacts \emph{improves} by $+24\%$.
\end{abstract}
